\section{The scope of the strengthened theory}\label{sec:current-scope}
The boundaries in Appendix~\ref{sec:boundaries} concern the inherited causal transfer theorem. The present synthesis theorems add necessary-and-sufficient conclusions on explicitly stated classes; they do not turn that sufficient transfer theorem into a classification of all standard Borel causal experiments.

The task-defect theorem is exact for a fixed selected loss family, its declared decision carrier, and one global finite decoder. Its occupation version holds for an actual occupation measure supplied by a competitor. Neither assertion manufactures that measure or an acquisition program. The scalar squared-task order is strictly weaker than the affine-square matrix order; the latter remains an exact characterization of its own larger test family. The Bregman comparison preserves squared-loss exponents only under the stated curvature bounds.

The finite-obstruction hierarchy characterizes one common full-history program for finitely branching instruments, bounded terminal scores and a strictly positive call-cost floor. It permits arbitrarily many operating models, but finite testing of an arbitrary uncountable family is not a decision procedure for its infinite synthesis problem. Compactness produces a policy, not a uniformly bounded number of live labels. Its finite dual concerns one specified horizon and finite model set. Fallback completion requires a genuinely legal common continuation and charges its cost and memory. The current supplement treats deadlocks, restricted preparation faces and the distinction between small survival mass and small expected cost tails.

The finite-memory certificate is intrinsic to a fixed finite-label observer on a declared time-homogeneous interface. Its projections and potential matrices are proof certificates, not stored observer variables. The finite algebraic decision conclusion requires finitely many models, algebraic primitive data and a semialgebraic decision carrier. It does not assert efficient synthesis, rational optimal probabilities, an exact finite-bit implementation of arbitrary algebraic randomization, or optimal program, scratch and arithmetic lower bounds. Exact initial-state uniformity does not replace a missing certificate at later roots of a history-dependent experiment. Returning from a balancing subroutine must be compiled at the current productive report before its prospective audit, as specified in the supplement.

The two-detector theorem has no balancing group and no free reset. It characterizes all integer classical live-memory budgets for two projectors, unit call costs, first-success commitment and no other physical command. The maximum principal-angle cosine is the worst spectral block of that model, not a universal rank or memory exponent. The initial-state saddle at unrestricted memory is proved from the physical instrument. The exact two-label minimax gap shows that a single preparation need not supply a sharp robust converse after imposing a nonconvex controller restriction. It does not imply such a gap for every memory-constrained experiment. Randomized initial labels are a separately stated interface, not an uncharged repair of a fixed-start observer.

All inherited valid results remain either in the native main inputs, the complete preceding companion, or their pinned independent realization lineages. In particular, the ordered-measurement local geometry and allocation theory, exact two-stage reset and message-discrimination theory, rank constraints and spectral variational values are not asserted to follow from the detector example. Conversely the current article does not abandon them because its general structural results have become central. The accompanying preservation and dependency maps identify their locations and prevent a change of research identity from being hidden in a reorganization.

Open mathematical goals remain: an intrinsic classification beyond finite branching; minimal memory under arbitrary resource-certified simulators; a broad nonintervening controlled-sensing application beyond the retained directional model; matching computational-resource lower bounds; and the independent unbounded-operator, large-deviation and wider A1/A2 pipeline targets. These are not hypotheses standing in for a missing proof of a theorem stated here. The present proofs establish their displayed conclusions with their displayed quantifiers. Finite diagnostics test particular algebraic identities, boundary cases and implementations; they are not formal proofs, independent expert assessment or an editorial decision on novelty and significance.
